
\section{\trSerreDuality{}的一瞥}

\trRiemannRoch{}的一个常见版本牵涉\trSerreDuality；后者与\trRiemannRoch{}不同，是难以证明的。

\textbf{18.5.1.\ \trhTheorem（光滑射影簇的\trSerreDuality）。}\ ——\ \emph{设 $X$ 是一个\trGeometricallyIrreducible\trSmooth\trProjective{} $k$-\trVariety，\trDimension{}为 $n$。那么 $X$ 上存在一个\trInvertibleSheaf{} $\omega_X$（或简记作 $\omega$），使得}
\[
h^i(X,\mathscr{F})=h^{n-i}(X,\omega_X\otimes\mathscr{F}^{\vee})
\]
\emph{对一切 $i\in\mathbb{Z}$ 与一切有限\trRank{}\trLocallyFreeSheaf{} $\mathscr{F}$ 成立。}

这个\trInvertibleSheaf{} $\omega_X$ 称作\emph{\trCanonical{}层}或\emph{\trDualizingSheaf}，它将在 \S29.1.5 中给出正式定义。我们将在第 29 章看到：\mbox{定理~18.5.1} 是下述\trPerfectPairing
\begin{equation*}\tag{18.5.1.1}\label{eq:serre-duality-perfect-pairing}
\mathrm{H}^i(X,\mathscr{F})\times\mathrm{H}^{n-i}(X,\omega_X\otimes\mathscr{F}^{\vee})\longrightarrow\mathrm{H}^n(X,\omega_X)\cong k,
\end{equation*}
的一个推论，并且光滑性还可以稍稍放宽。

\textbf{18.5.2.}\ \emph{进一步的奇迹：代数体积形式层是\trSerre{}对偶化的。}\quad 这个\trInvertibleSheaf{} $\omega_X$ 结果竟是 $X$ 的\trCotangentBundle{} $\Omega_{X/k}$ 的\trDeterminant{} $\mathscr{K}_X$（见 \S21.5.3）；我们将在第 21 章中定义\trCotangentBundle。我们将在 \S29.4 中建立 $\omega_X=\mathscr{K}_X$ 这一联系；参看\mbox{愿望~29.1.1}（Desideratum）。

\textbf{18.5.3.}\ \emph{回到\trRiemannRoch。}\quad 为了把曲线 $C$ 的\trRiemannRoch{}加以重述，只须注意到 $h^1(C,\mathscr{L})=h^0(C,\omega_C\otimes\mathscr{L}^{\vee})$ 便已足够。于是\trRiemannRoch{}公式可以改写作
\[
h^0(C,\mathscr{L})-h^0(C,\omega_C\otimes\mathscr{L}^{\vee})=\deg\mathscr{L}-p_a(C)+1.
\]

\textbf{18.5.A.（练习；假定\trSerreDuality）}\quad 设 $C$ 是 $k$ 上一条\trGeometricallyIntegral\trSmooth\trProjective{}\trCurve。

（a）证明：$h^0(C,\omega_C)$ 就是 $C$ 的\trGenus{} $g$。（亏格 $g$ 曾以\trArithmeticGenus{} $p_a$ 的面目定义于 \S18.4.4--18.4.5。）

（b）证明：$\deg\omega_C=2g-2$。（提示：取 $\mathscr{L}=\omega_C$ 的\trRiemannRoch。）

\textbf{18.5.4.}\ \emph{\trhExample{。}}\quad 若 $C=\mathbb{P}^1_k$，则\mbox{练习~18.5.A} 蕴涵 $\omega_C\cong\OO(-2)$。此外，我们还有一个自然的\trPerfectPairing{}（对照 \eqref{eq:serre-duality-perfect-pairing}）
\[
\mathrm{H}^0(\mathbb{P}^1,\OO(n))\times\mathrm{H}^1(\mathbb{P}^1,\OO(-2-n))\longrightarrow k\cong\mathrm{H}^1(\mathbb{P}^1,\OO(-2)).
\]
我们可以按下面的方式来解读这个配对。若 $n<0$，则左边的两个因子都是 $0$，故设 $n>0$。此时 $\mathrm{H}^0(\mathbb{P}^1,\OO(n))$ 对应于 $x$、$y$ 的 $n$ 次\trHomogeneous{}多项式，而 $\mathrm{H}^1(\mathbb{P}^1,\OO(-2-n))$ 对应于 $x$、$y$ 的 $-2-n$ 次\trHomogeneous{} Laurent 多项式，其中 $x$ 与 $y$ 的次数都不超过 $-1$（见\mbox{注记~18.3.5}）。你可以很快地验证：这两个向量空间的维数都是 $n+1$。这个配对给出如下：把多项式与那个 Laurent 多项式相乘，便得到一个次数为 $-2$ 的 Laurent 多项式；再读出 $x^{-1}y^{-1}$ 的系数。（这套做法对 $\mathbb{P}^n_A$ 也行之有效；见\mbox{定理~18.1.2} 陈述之后的讨论。）

\textbf{18.5.B.（练习；假定\trSerreDuality：连通光滑射影簇上的\trAmple\trDivisor{}是连通的）}\quad 设 $X$ 是一个\trDimension{}至少为 $2$ 的\trConnected\trSmooth\trProjective{} $k$-\trVariety，而 $D$ 是一个\trEffectiveDivisor，且是\trAmple{}的。证明：$D$ 是\trConnected{}的。提示：设 $D=V(s)$，其中 $s$ 是某个\trAmple\trInvertibleSheaf{} $\mathscr{L}$ 的一个\trSection。那么对一切 $n>0$ 都有 $V(s^n)=V(s)$（作为\emph{\trVanishingSet}而非\emph{\trVanishingScheme}），故我们可以把 $\mathscr{L}$ 换成由我们选定的高次幂。利用 $0\to\OO_X(-nD)\to\OO_X\to\OO_{V(s^n)}\to 0$ 的\trLongExactSequence，证明：当 $n\gg 0$ 时，$h^0(X,\OO_{V(s^n)})=1$。（这或许会让你想起 \S11.4.7。）

一旦我们得知\trSerreDuality{}对\trCM{}射影概形也成立（\S29.3），这个结果便会自动地推广到这些概形上。（一个有关的结果是\mbox{练习~18.6.U}，它并不使用\trSerreDuality。）另一方面，若 $X$ 是 $\mathbb{P}^4$ 中交于一点的两个二维平面之并，则这个结果不真（为什么？），于是这将蕴含：这个 $X$ 不是\trCM{}的。（我们将在\mbox{反例~26.2.2} 中以另一种方式证明这一点。）

\textbf{18.5.5.}\ \hard\ \textbf{$\mathbb{P}^1_k$ 上\trVectorBundle{}的分类。}

正如\mbox{例~17.2.6} 中所许诺的，我们将对 $\mathbb{P}^1_k$ 上的\trVectorBundle{}进行分类。我们把这一讨论置于 \S18.5，因为你在其中一步或许会用到\trSerreDuality（尽管你也可能用不上）。

\textbf{18.5.6.\ \trhTheorem{。}}\ ——\ \emph{若 $\mathscr{E}$ 是 $\mathbb{P}^1_k$ 上\trRank{}为 $r$ 的\trVectorBundle，则 $\mathscr{E}\cong\OO(a_1)\oplus\cdots\oplus\OO(a_r)$，其中 $a_1,\dots,a_r$ 是一个唯一的非降整数序列。}

这个结果曾被独立地证明过许多次（例如见 \trDedekind{}与\trWeber{}的文章 [DW, \S22]），而且它是\trGrothendieck{}的一个定理 [Gr2] 的特例。它有时被称作\trGrothendieck{}定理——因为以\trGrothendieck{}命名的定理还不够多。

\textbf{18.5.7.}\quad 对 $\mathbb{P}^n_k$ 更一般地而言，$r=1$ 的情形已在 \S15.4.15 中证明，但当 $r>1$ 且 $n>1$ 时该命题不真。\mbox{练习~21.3.H} 给出了一个反例，它表明：甚至连用\trLineBundle{}滤过 $\mathbb{P}^2$ 上\trRank{}为 $2$ 的丛 $\Omega_{\mathbb{P}^2}$ 都做不到。一个确实成立的推广是 Horrocks 的一个定理：$\mathbb{P}^n_k$ 上的有限\trRank{}\trLocallyFreeSheaf{} $\mathscr{E}$ 分裂，恰好当「它的一切扭转的中间\trCohomology{}全为零」——也就是说，当对 $0<i<n$ 与一切 $m\in\mathbb{Z}$ 都有 $\mathrm{H}^i(\mathbb{P}^n_k,\mathscr{E}(m))=0$。这有一个出人意料的推论：若 $n\ge 2$，则 $\mathbb{P}^n_k$ 上的一个有限\trRank{}\trLocallyFreeSheaf{}分裂，当且仅当它限制到某张预先取定的二维平面（$\mathbb{P}^2_k$）上时分裂——一切额外的复杂性都已经在维数 $2$ 处现身。更多内容见 [OSS, \S2.3]。

\emph{证明。}\quad 注意，这个分类对\trCohomology{}只字未提，但其证明却在本质之处用到了\trCohomology。证明\mbox{定理~18.5.6} 时，可以不用任何上同调机制（例如见 [HM] 或 [GW, p.~314-5]），也可以动用更多的上同调机制（例如见 [Ha1, Ex.~V.2.6]）。

\textbf{18.5.C.（练习）}\quad 设 $\mathscr{E}\cong\OO(a_1)\oplus\cdots\oplus\OO(a_r)$，其中 $a_1\le\cdots\le a_r$。证明：诸 $a_i$ 可以由 $\dim_k\Hom(\OO(m),\mathscr{E})$ 这些数（$m$ 遍历诸整数）确定。用这一点证明\mbox{定理~18.5.6} 的唯一性部分。

现在我们开始证明，对 $r$ 作归纳。取定一个\trRank{}为 $r$ 的\trLocallyFreeSheaf{} $\mathscr{E}$。$r=1$ 的情形已在 \S15.4.15 中确立，故设 $r>1$。

\textbf{18.5.D.（练习）}\quad 证明：当 $m\ll 0$ 时，$\Hom(\OO(m),\mathscr{E})\ne 0$；而当 $m\gg 0$ 时，$\Hom(\OO(m),\mathscr{E})=0$。提示：用\mbox{练习~14.1.G} 证明 $\Hom(\OO(m),\mathscr{E})=\mathscr{E}(-m)$。第一部分使用\trSerre{}定理 A（\mbox{定理~16.1.5}）。第二部分请随意使用\trSerreDuality{}与\trSerre{}消失。不过你或许更愿意想出一个不用\trSerreDuality{}的论证，以免援引我们尚未证明的东西。

于是存在某个 $a_r$，使得 $\dim\Hom(\OO(a_r),\mathscr{E})>0$，而对一切 $m>a_r$ 都有 $\Hom(\OO(m),\mathscr{E})=0$。选取一个非零映射 $\phi\colon\OO(a_r)\to\mathscr{E}$。

\textbf{18.5.E.（练习）}\quad 证明：$\phi$ 是\trInjectiveMap。（提示：$\OO(a_r)$ 是\trTorsionFree{}的，从而其\trKernel{}也是\trTorsionFree{}的。）

\textbf{18.5.F.（练习）}\quad 设 $\mathscr{F}$ 是 $\phi$ 的\trCokernel。证明：$\mathscr{F}$ 是\trLocallyFree{}的。提示：\mbox{练习~14.3.F(c)} 给出了一个\trExactSequence{} $0\to\mathscr{F}_{\mathrm{tors}}\to\mathscr{F}\to\mathscr{F}_{\mathrm{lf}}\to 0$，其中 $\mathscr{F}_{\mathrm{lf}}$ 是\trLocallyFree{}的。令 $\mathscr{L}$ 为满射 $\mathscr{E}\to\mathscr{F}_{\mathrm{lf}}$ 的\trKernel。证明：$\mathscr{L}$ 是\trLocallyFree{}的，从而同构于某个 $\OO(N)$。证明：存在非零映射 $\phi'\colon\OO(a_r)\to\OO(N)$。证明（用与上一道练习相同的想法）：这个非零映射 $\phi'$ 必是\trInjectiveMap，于是 $a_r\le N$。证明 $N\le a_r$，因为存在非零映射 $\OO(N)\to\mathscr{E}$（回想一下 $a_r$ 当初是怎么选的）。证明：$\phi'$ 是\trIsomorphism，从而 $\mathscr{F}=\mathscr{F}_{\mathrm{lf}}$。

由归纳假设，我们有 $\mathscr{F}=\OO(a_1)\oplus\cdots\oplus\OO(a_{r-1})$，其中 $a_1\le\cdots\le a_{r-1}$，于是我们有一条\trShortExactSequence
\begin{equation*}\tag{18.5.7.1}\label{eq:p1-bundles-ses}
0\longrightarrow\OO(a_r)\longrightarrow\mathscr{E}\longrightarrow\OO(a_1)\oplus\cdots\oplus\OO(a_{r-1})\longrightarrow 0.
\end{equation*}

接下来我们证明：对 $i<r$ 有 $a_r\ge a_i$。我们的动机来自这样一个事实：对 $\mathbb{P}^1$ 上的\trQCSheaf{} $\mathscr{G}$，
\[
\Hom_{\mathbb{P}^1}(\OO(a_r+1),\mathscr{G})=\mathrm{H}^0(\mathbb{P}^1,\mathscr{G}(-a_r-1))
\]
（\mbox{练习~14.1.G}）。用 $\OO(-a_r-1)$ 去张量 \eqref{eq:p1-bundles-ses}（由\mbox{练习~14.1.F}，正合性得以保持），并取\trCohomology{}中的\trLongExactSequence。该\trLongExactSequence{}的一部分是
\[
\mathrm{H}^0(\mathbb{P}^1,\mathscr{E}(-a_r-1))\longrightarrow\mathrm{H}^0(\mathbb{P}^1,\mathscr{F}(-a_r-1))\longrightarrow\mathrm{H}^1(\mathbb{P}^1,\OO(-1)).
\]

注意 $\mathrm{H}^0(\mathbb{P}^1,\mathscr{E}(-a_r-1))=0$（因为由 $a_r$ 的定义，$\Hom(\OO(a_r+1),\mathscr{E})=0$），且 $\mathrm{H}^1(\mathbb{P}^1,\OO(-1))=0$（\mbox{定理~18.1.2(b)}）。于是
\[
\begin{gathered}
0=\mathrm{H}^0(\mathbb{P}^1,\mathscr{F}(-a_r-1))\\
=\mathrm{H}^0(\mathbb{P}^1,\oplus_i\OO(a_i-a_r-1))\\
=\oplus_i\mathrm{H}^0(\mathbb{P}^1,\OO(a_i-a_r-1)).
\end{gathered}
\]
于是 $a_i-a_r-1<0$（由\mbox{练习~15.1.A}），故 $a_r\ge a_i$，正如所愿。

最后，我们想证明：正合列 \eqref{eq:p1-bundles-ses} 把 $\mathscr{E}$ 表示成了一个\trDirectSum（即\trSubsheaf{}与\trQuotientSheaf{}的直和）。为简单起见，我们只关注 $r=2$ 的情形，一般情形则留待\mbox{练习~18.5.I}。

\textbf{18.5.G.（练习）}\quad 证明：在适当的\trCoordinate{}下，这个\trVectorBundle{}的\trTransitionFunction{}由
\[
\begin{pmatrix} t^{-a_2} & \alpha(t) \\ 0 & t^{-a_1} \end{pmatrix},
\]
给出，其中 $\alpha(t)$ 是 $t$ 的一个 Laurent 多项式。提示／提醒：回想一下，$\mathbb{A}^1_k$ 上的一切\trVectorBundle{}都是\trTrivial{}的（\mbox{练习~14.2.C}）。关于以一个\trVectorBundle{}（转移矩阵已知）扩张另一个\trVectorBundle{}的转移矩阵，曾在\mbox{练习~14.2.G} 中讨论过。

\textbf{18.5.H.（练习）}\quad 隐含在上述 $2\times 2$ 矩阵之中的，是如下两组基的选取：环 $k[t]$ 上一个\trRank{}为 $2$ 的\trFreeModule{} $M$ 的一组基，相应于诸标准\trAffine{}开子集 $\Spec k[t]$ 中的某一个；以及在另一张标准开子集之上的、环 $k[1/t]$ 上一个\trRank{}为 $2$ 的\trFreeModule{} $M'$ 的一组基。证明：对 $M$ 的基施以一个适当的「上三角」更换，便可使 $\alpha(t)$ 不含次数 $\ge -a_2$ 的\trMonomial。证明：对 $M'$ 的基施以一个适当的「上三角」更换，便可使 $\alpha(t)$ 不含次数 $\le -a_1$ 的\trMonomial。于是，适当选取 $M$ 与 $M'$ 的基之后，我们便可取 $\alpha(t)=0$。证明：这蕴含 $\mathscr{E}=\OO(a_1)\oplus\OO(a_2)$。

\textbf{18.5.I.（练习）}\quad 对任意的 $r$ 完成\mbox{定理~18.5.6} 的证明。提示：在 $r=2$ 的情形之外，再没有任何新想法。\hfill$\square$
